\documentclass{beamer}

\usepackage[onehalfspacing]{setspace} %space between lines
\usepackage[utf8]{inputenc}
\usepackage{gensymb} % for the degree symbol
\usepackage{url}
\usepackage{mathtools} % multi-choice equations

\usepackage[
backend=biber,
style=numeric,
sorting=ynt,
maxbibnames=12 
]{biblatex}
%\addbibresource{Sources.bib}


\usetheme{Madrid}
\usecolortheme{default}

%------------------------------------------------------------
%This block of code defines the information to appear in the
%Title page
\title[Mathematical Optimization Algorithms] % bottom middle
{Spaceship Titanic}

\subtitle{Mathematical Optimization Algorithms}

\author[] % (optional)
{}

\institute[] % (optional)
{
%   \inst{1}%
%   Faculty of Mathematics and Computer Science\\
%   Universität Heidelberg
%   \and
%   \inst{2}%
%   Faculty of ??\\
%   Politechnika Warszawska
}

\date[WS 2025/2026] % (optional)
{Winter Semester 2025/2026, January 2026}



%End of title page configuration block
%------------------------------------------------------------



%------------------------------------------------------------
%The next block of commands puts the table of contents at the 
%beginning of each section and highlights the current section:

\AtBeginSection[]
{
\begin{frame}
    \frametitle{Table of Contents}
    \tableofcontents[currentsection]
  \end{frame}
}
%------------------------------------------------------------


\begin{document}

%The next statement creates the title page.
\frame{\titlepage}


%---------------------------------------------------------
%This block of code is for the table of contents after
%the title page
\begin{frame}
\frametitle{Table of Contents}
\tableofcontents
\end{frame}
%---------------------------------------------------------

\section{Introduction}
\begin{frame}{Introduction}
    \begin{itemize}
        \item 
        \item
    \end{itemize}
\end{frame}

%---------------------------------------------------------
\begin{frame}{Goals}
    \begin{itemize}
        \item Accuracy: 95\% or better
        \item Trying different: 
            \begin{itemize}
                \item layers
                \item neuron counts
                \item optimizers
                \item hyperparameters (such as learning rate or weight decay)
            \end{itemize}
    \end{itemize}
\end{frame}

%---------------------------------------------------------

\begin{frame}{Data}
    \begin{itemize}
        \item training file, 8694 rows, column names:
            \begin{itemize}
                \item PassengerId
                \item HomePlanet
                \item CryoSleep
                \item Cabin
                \item Destination
                \item Age
                \item VIP
                \item RoomService, FoodCourt, ShoppingMall, Spa, VRDeck
                \item Name
                \item Transported
            \end{itemize}
        \item test file, 4277 rows, same columns as training but without 'Transported'
    \end{itemize}
\end{frame}


\begin{frame}{Data}
    \begin{itemize}
        \item training file, 8694 rows, column names:
            \begin{itemize}
                \item PassengerId
                \item HomePlanet
                \item CryoSleep
                \item Cabin
                \item Destination
                \item Age
                \item VIP
                \item RoomService, FoodCourt, ShoppingMall, Spa, VRDeck
                \item Name
                \item Transported
            \end{itemize}
        \item test file, 4277 rows, same columns as training but without 'Transported'
    \end{itemize}
\end{frame}
%---------------------------------------------------------
\section{Code}
\begin{frame}{Preprocessing}
Some values are not numerical
\begin{itemize}
    \item 
    \item
\end{itemize}

\end{frame}

\begin{frame}{Neural network architecture}
    \begin{itemize}
        \item Normalization
        \item Dropout (in layers %TODO
        \item Layers:
        \item Optimizers:
        \item Activation functions: 
    \end{itemize}
    
\end{frame}


\section{Experimentation}
\begin{frame}{Process}
    %Todo explain how we proceeded
\end{frame}

\begin{frame}{Experimentation}
    
\end{frame}

\begin{frame}{Conclusions/Issues}
    
    
\end{frame}

\begin{frame}{Sources}

%\nocite{*}
%\printbibliography
    

\end{frame}




\end{document}
